%=======================================================================
% anexoA.tex — Conjunto de especies objetivo
%=======================================================================
La \tabref{tab:anexo-especies} lista las 55 especies del conjunto objetivo con
su categoría vigente en el Reglamento de Clasificación de Especies
\citep{mma2026rce}, su condición de endémica, el grupo al que pertenecen dentro
del diseño del corpus y la disponibilidad de datos según el censo de metadatos
realizado en Xeno-canto y GBIF. ``NE'' indica que la especie no ha sido
evaluada por el RCE. La columna ``XC Chile'' es la que define el corpus de
audio de este trabajo; ``XC global'' se reporta solo para las especies de
interés de conservación, para las que se relevó la disponibilidad fuera del
país.

\setlength{\tabcolsep}{2.5pt}
{\scriptsize
\begin{longtable}{@{}l l c c l r r r@{}}
  \caption{Conjunto objetivo de 55 especies y disponibilidad de datos.}
  \label{tab:anexo-especies}\\
  \toprule
  Nombre científico & Nombre común & RCE & Endémica & Grupo & XC Chile & XC global & Fotos GBIF \\
  \midrule
  \endfirsthead
  \multicolumn{8}{l}{\footnotesize\textit{Tabla A.1 (continuación)}}\\
  \toprule
  Nombre científico & Nombre común & RCE & Endémica & Grupo & XC Chile & XC global & Fotos GBIF \\
  \midrule
  \endhead
  \bottomrule
  \endlastfoot
    \textit{Laterallus jamaicensis} & pidencito & EN & No & Amenazada & 34 & 181 & 42 \\
    \textit{Conirostrum tamarugense} & comesebo de los tamarugales & EN & No & Amenazada & 22 & 23 & 38 \\
    \textit{Chloephaga melanoptera} & piuquén & EN & No & Amenazada & 17 & 55 & 421 \\
    \textit{Rallus antarcticus} & pidén austral & VU & No & Amenazada & 16 & 22 & 24 \\
    \textit{Chloephaga picta} & caiquén & VU & No & Amenazada & 14 & 38 & 991 \\
    \textit{Cyanoliseus patagonus} & tricahue & EN & No & Amenazada & 14 & 55 & 414 \\
    \textit{Recurvirostra andina} & caití & VU & No & Amenazada & 6 & 32 & 130 \\
    \textit{Limosa haemastica} & zarapito de pico recto & VU & No & Amenazada & 4 & 31 & 261 \\
    \textit{Agriornis albicauda} & mero de la puna & VU & No & Amenazada & 2 & 23 & 5 \\
    \textit{Anous stolidus} & gaviotín de San Félix & VU & No & Amenazada & 0 & 41 & 11 \\
    \textit{Calidris canutus} & playero ártico & EN & No & Amenazada & 0 & 311 & 21 \\
    \textit{Charadrius nivosus} & chorlo nevado & VU & No & Amenazada & 0 & 39 & 2 \\
    \textit{Podiceps gallardoi} & pimpollo tobiano & CR & No & Amenazada & 0 & 40 & 0 \\
    \textit{Pteroptochos megapodius} & turca & NE & Sí & Endémica & 86 &  &  \\
    \textit{Scelorchilus albicollis} & tapaculo & NE & Sí & Endémica & 65 &  &  \\
    \textit{Mimus thenca} & tenca & NE & Sí & Endémica & 48 &  &  \\
    \textit{Scytalopus fuscus} & churrín del norte & NE & Sí & Endémica & 46 &  &  \\
    \textit{Pteroptochos castaneus} & hued-hued castaño & NE & Sí & Endémica & 41 &  &  \\
    \textit{Enicognathus leptorhynchus} & choroy & LC & Sí & Endémica & 35 & 35 & 258 \\
    \textit{Ochetorhynchus melanurus} & chiricoca & NE & Sí & Endémica & 34 &  &  \\
    \textit{Turdus falcklandii} & zorzal & NE & No & Común & 186 &  &  \\
    \textit{Aphrastura spinicauda} & rayadito & NE & No & Común & 105 &  &  \\
    \textit{Troglodytes aedon} & chercán & NE & No & Común & 102 &  &  \\
    \textit{Zonotrichia capensis} & chincol & NE & No & Común & 92 &  &  \\
    \textit{Scelorchilus rubecula} & chucao & LC & No & Común & 92 &  &  \\
    \textit{Curaeus curaeus} & tordo & NE & No & Común & 88 &  &  \\
    \textit{Elaenia chilensis} & fío-fío & NE & No & Común & 85 &  &  \\
    \textit{Eugralla paradoxa} & churrín de la Mocha & NE & No & Común & 84 &  &  \\
    \textit{Scytalopus magellanicus} & churrín del sur & NE & No & Común & 75 &  &  \\
    \textit{Leptasthenura aegithaloides} & tijeral & NE & No & Común & 68 &  &  \\
    \textit{Sephanoides sephaniodes} & picaflor chico & NE & No & Común & 67 &  &  \\
    \textit{Pseudasthenes humicola} & canastero & NE & No & Común & 67 &  &  \\
    \textit{Anairetes parulus} & cachudito & NE & No & Común & 56 &  &  \\
    \textit{Sylviorthorhynchus desmurii} & colilarga & NE & No & Común & 55 &  &  \\
    \textit{Glaucidium nana} & chuncho & NE & No & Común & 54 &  &  \\
    \textit{Spinus barbatus} & jilguero & NE & No & Común & 54 &  &  \\
    \textit{Pteroptochos tarnii} & hued-hued del sur & LC & No & Común & 52 &  &  \\
    \textit{Phrygilus patagonicus} & cometocino patagónico & NE & No & Común & 51 &  &  \\
    \textit{Phrygilus gayi} & cometocino de Gay & NE & No & Común & 46 &  &  \\
    \textit{Pardirallus sanguinolentus} & pidén & NE & No & Común & 45 &  &  \\
    \textit{Vanellus chilensis} & queltehue & NE & No & Común & 45 &  &  \\
    \textit{Geositta rufipennis} & minero cordillerano & NE & No & Común & 44 &  &  \\
    \textit{Anthus correndera} & bailarín chico & NE & No & Común & 43 &  &  \\
    \textit{Veniliornis lignarius} & carpinterito & NE & No & Común & 42 &  &  \\
    \textit{Leistes loyca} & loica & NE & No & Común & 41 &  &  \\
    \textit{Diuca diuca} & diuca & NE & No & Común & 41 &  &  \\
    \textit{Strix rufipes} & concón & NT & No & Común & 39 &  &  \\
    \textit{Rhopospina fruticeti} & yal & NE & No & Común & 38 &  &  \\
    \textit{Pygarrhichas albogularis} & comesebo grande & NE & No & Común & 38 &  &  \\
    \textit{Sicalis luteola} & chirihue & NE & No & Común & 38 &  &  \\
    \textit{Pyrope pyrope} & diucón & NE & No & Común & 37 &  &  \\
    \textit{Cinclodes fuscus} & churrete acanelado & NE & No & Común & 36 &  &  \\
    \textit{Cinclodes patagonicus} & churrete & NE & No & Común & 33 &  &  \\
    \textit{Colorhamphus parvirostris} & viudita & NE & No & Común & 32 &  &  \\
    \textit{Tachycineta leucopyga} & golondrina chilena & NE & No & Común & 32 &  &  \\
\end{longtable}}
